% The brief: the results, in the order a paper would present them, each with the
% certificate that backs it.  This is the source the human-written exposition in expos/ draws
% on; that document may reorder and reword freely, but where the two disagree about a number
% or a claim, this one is right.
%
% It is formatted as an appendix so that the body's equation numbering does not shift.  The
% labels are stable and nothing in the body cites them.

\appendix
% article's \appendix renumbers sections but not equations, so the counter is reset here and
% given an A prefix.
\setcounter{equation}{0}
\renewcommand{\theequation}{A\arabic{equation}}
\section{Not Really an Appendix}
\label{sec:summary}

This is the whole calculation, stated and not derived: the results, in the order a paper
would present them, each beside the certificate that backs it.  The sections before this one
are the derivations.

Conventions throughout: the sky axes are $x$ towards North and $y$ towards East, the third
axis points at the observer, and a map is normalised so that its monopole is
$a_{00} = 1$, which is the sphere mean itself.  The star has unit angular radius, so the
visible disk is the unit circle.

\paragraph{What the certificates are.}  Each is a function in
\code{mostlyharmless.symbolic.derivations}, called by a one-line assertion in
\code{make check}.  Each is a \emph{differential} test: two independent computations of the
same number, compared to a stated tolerance, with the residual returned so that degradation is
visible.  They fall into three classes.

\begin{description}
\item[Theorem-like] the check states a result the reader may want, and its reference is the
  definition or a route no other check shares.  These are worth reading.
\item[Lemma-like] a subsidiary identity the derivation uses.  Short to state, and where a
  convention error hides.
\item[Scaffolding] a check that a later result makes redundant.  Useful while the
  derivation was being found; skippable once the terminal check has passed.
\end{description}

Two limits are worth stating plainly.  A certificate is only as good as its independence: where
both routes share a step, a wrong step passes, and V24 says so of itself.  And the trusted base
is numpy, scipy and sympy, not a proof kernel --- so a reader who wants to check the checking
must read the script, which is why each is ten to thirty lines.

\subsection{The visibility}

The complex visibility is the Fourier transform of the sky brightness with $2\pi$ inside the
exponent and no prefactor,
\begin{equation}
  V(u,v) \;=\; \int I(x,y)\, e^{-2\pi i (u x + v y)}\, dx\, dy .
  \label{eq:brief-visibility}
\end{equation}
It is left unnormalised, so $V(0,0)$ is the total flux; the normalisation lives in the
$\Delta V$ the models plot.

The transform is done in plane polar coordinates on the sky and on the baseline plane.  The
azimuth $\varphi$ runs from North through East on the sky and the position angle $\psi$ does the
same on the ground, so $x$ and $u$ are the North components of both; $\rho$ is the projected
radius, so the visible disk is $\rho \le 1$, and $q$ is the baseline length:
\begin{align}
  x = \rho\cos\varphi , \qquad y = \rho\sin\varphi ,
    \label{eq:brief-sky-polar} \\
  u = q\cos\psi , \qquad v = q\sin\psi .
    \label{eq:brief-baseline-polar}
\end{align}

\begin{itemize}
  \item[$\checkmark$] \code{check_uniform_disk_visibility} checks the uniform disk against
    $2J_1(x)/x$, with flux $\pi$, the disk's area --- the anchor that fixes the harmonic
    normalisation, the placement of $2\pi$ and the planar measure jointly.
    \\[1pt] \vtag{16}\quad \emph{Numerical, $l = 0$}
\end{itemize}

\subsection{Spherical harmonics}

The map is carried in complex harmonics,
\begin{equation}
  Y_{lm}(\theta,\varphi) = N_{lm}\, P_{lm}(\cos\theta)\, e^{im\varphi},
  \qquad
  N_{lm} = \sqrt{(2l+1)\,\frac{(l-m)!}{(l+m)!}} ,
  \label{eq:brief-ylm}
\end{equation}
The associated functions come from Rodrigues' formula for the Legendre polynomials,
\begin{equation}
  P_{lm}(x) \;=\; (1-x^2)^{m/2}\,\frac{d^m P_l}{dx^m} .
  \label{eq:brief-rodrigues}
\end{equation}
The square root in $N_{lm}$ is taken positive, so the only arbitrary sign left is the one in
this display.  The Condon--Shortley convention fixes it as $(-1)^m$; this work does not, taking
\eqref{eq:brief-rodrigues} as it stands, which is why no $(-1)^m$ appears in the closed forms
below.  Only $m \ge 0$ is ever needed: reality is a property of the expansion rather than of a
separate set of coefficients, so half the orders suffice and the signed order is never formed.

The normalisation is $Y_{00} = 1$, so the average of $|Y_{lm}|^2$ over the sphere is unity and
its integral is $4\pi$; the orthonormal convention scales every $Y_{lm}$ by $1/\sqrt{4\pi}$
relative to this one.

At $l = 2$ the definition gives three stored values in closed form, and they are what the
deformation uses: \eqref{eq:brief-zonal} combines them into a single zonal pattern,
\begin{equation}
  \begin{aligned}
    Y_{20} &= \sqrt{\tfrac54}\,\bigl(3\cos^2\theta - 1\bigr), \\
    Y_{21} &= \sqrt{\tfrac{15}{2}}\,\cos\theta\,\sin\theta\; e^{i\varphi}, \\
    Y_{22} &= \sqrt{\tfrac{15}{8}}\,\sin^2\theta\; e^{2i\varphi} .
  \end{aligned}
  \label{eq:brief-y2m}
\end{equation}

\begin{itemize}
  \item[$\checkmark$] \code{check_legendre_orthogonality_exact} integrates $P_lP_{l'}$ in rational
    arithmetic.
    \\[1pt] \vtag{1}\quad \emph{Exact, $l \le 5$}
  \item[$\checkmark$] \code{check_reality_exact} checks that the paired coefficients of a real
    map satisfy their reality condition identically.
    \\[1pt] \vtag{2}\quad \emph{Exact}
  \item[$\checkmark$] \code{check_real_basis_orthonormality_numeric} checks the real basis on the
    grid, which is what makes it usable for rendering.
    \\[1pt] \vtag{25}\quad \emph{Numerical, $l \le 4$}
\end{itemize}

\subsection{Rotation}

A map carried in the body frame reaches the observer through a rotation of its coefficients,
$b = D(R)a$, where $D$ is the Wigner matrix of the rotation $R$.  In the zyz convention
$R = R_z(\gamma)R_y(\beta)R_z(\alpha)$, with $\alpha$ the spin phase, $\beta$ the inclination
and $\gamma$ the sky position angle of the spin.  Two features do the work: it is block
diagonal in $l$, so the degree is untouched and a truncated map stays truncated; and each block
is a finite matrix of $(2l+1)^2$ numbers depending only on the three angles, so a rotation costs
a matrix product rather than a resampling.

For the zonal $l = 2$ patterns that the deformation needs, this specialises to one identity.
With $\hat e$ the axis of the pattern,
\begin{equation}
  P_2(\hat n\cdot\hat e) \;=\; \frac{1}{5}\sum_{m=-2}^{2}
  \overline{Y_{2m}(\hat e)}\,Y_{2m}(\hat n) ,
  \label{eq:brief-zonal}
\end{equation}
which needs only the stored orders,\footnote{For an axis at direction cosines $(a,b,c)$ the
three stored values are \mbox{$\sqrt{5/4}(3c^2-1)$}, \mbox{$-\sqrt{15/2}\,c(a+ib)$} and
\mbox{$\sqrt{15/8}(a+ib)^2$}, which is the form needed when the axis is given as a vector
rather than as a pair of angles.}
so a zonal pattern costs three numbers in any frame.


\begin{itemize}
  \item[$\checkmark$] \code{check_addition_theorem_numeric} checks the addition theorem that
    \eqref{eq:brief-zonal} rests on, in its complex and its real form.
    \\[1pt] \vtag{3}\quad \emph{Numerical to $40$ digits, $l \le 6$}
  \item[$\checkmark$] \code{check_wigner_d_against_exponential} compares the $d$-matrix with
    $\exp(-i\beta J_y)$.
    \\[1pt] \vtag{4}\quad \emph{Numerical, $l \le 6$}
  \item[$\checkmark$] \code{check_wigner_d_properties} checks orthogonality and the symmetry of
    $d$ in its indices.
    \\[1pt] \vtag{5}\quad \emph{Numerical, $l \le 6$}
  \item[$\checkmark$] \code{check_rotation_against_geometry} and
    \code{check_rotation_nonnegative_form} check that a rotation maps the polar axis where the
    geometry says, and that the non-negative-order form agrees.
    \\[1pt] \vtag{6}\quad \emph{Numerical, $l \le 4$}
  \item[$\checkmark$] \code{check_rotation_composition} checks that two successive rotations
    compose as their matrices do.
    \\[1pt] \vtag{7}\quad \emph{Numerical, $l \le 3$}
\end{itemize}

\subsection{Combining the map with the transform}

A real map is expanded in complex harmonics,
\begin{equation}
  I(\theta,\varphi) \;=\; \sum_{l}\sum_{m=-l}^{l} a_{lm}Y_{lm} .
  \label{eq:brief-expansion}
\end{equation}
Written over the stored half, the same pairing works through the transform, and the parity of
$m$ decides which part of a coefficient survives: the contribution of one stored $a_{lm}$ is
twice the real part of $a_{lm}V_{lm}$ when $m$ is even and twice $i$ times its imaginary part
when $m$ is odd.  That is the whole of the negative-order bookkeeping, and it is why half the
coefficients suffice.  Its one consequence worth stating early: the two extensions differ, the
map's being a plain conjugate and the visibility's carrying $(-1)^m$, because reversing the
baseline advances the mode's azimuth by $m\pi$.


\begin{itemize}
  \item[$\checkmark$] \code{check_full_map_fft} compares the paired mode sum of
    \eqref{eq:brief-expansion} against a two-dimensional FFT of the same map, in the
    low-frequency band.
    \\[1pt] \vtag{15}\quad \emph{Numerical to $5\times10^{-2}$}
\end{itemize}

Because the disk is a circle, the angular part of \eqref{eq:brief-visibility} can be done
once and for all: each harmonic contributes a complex constant times a radial integral that
depends on the baseline length $q$ alone,
\begin{equation}
  V_{lm}(u,v) \;=\; 2\pi(-i)^m N_{lm}\, e^{im\psi}\, X_{lm}(q) ,
  \label{eq:brief-mode}
\end{equation}
with $\psi$ the baseline's position angle, $N_{lm}$ the harmonic normalisation, and
\begin{equation}
  X_{lm}(q) \;=\; \int_0^1 P_{lm}\bigl(\sqrt{1-\rho^2}\bigr)\, J_m(2\pi q\rho)\,
  \rho\, d\rho .
  \label{eq:brief-radial-integral}
\end{equation}
Everything about the geometry of a single mode is in $X_{lm}$; everything about its place on
the sky is in $e^{im\psi}$.  The reduction is exact at every $q$, and it is the object the
next two subsections evaluate in closed form.


\begin{itemize}
  \item[$\checkmark$] \code{check_jacobi_anger} checks the Jacobi--Anger expansion the reduction
    uses.
    \\[1pt] \vtag{17}\quad \emph{Numerical, at roundoff}
  \item[$\checkmark$] \code{check_baseline_parity} checks that reversing the baseline gives the
    conjugate, which is the parity the paired sum assumes.
    \\[1pt] \vtag{23}\quad \emph{Algebraic}
\end{itemize}

\subsection{The closed forms}

The integral splits by the parity of $l+m$, and the two halves are different in kind.

For $l+m$ odd, $P_{lm}(\sqrt{1-\rho^2})$ is $\rho^m\sqrt{1-\rho^2}$ times a polynomial in
$\rho^2$ of degree $(l-m-1)/2$; the factor of $\sqrt{1-\rho^2}$ is the one the Legendre
recurrence of section~\ref{sec:odd-recurrence} absorbs, and the Bessel integral then collapses
to a single term:
\begin{equation}
  X_{lm}(q) \;=\; \frac{(l+m)!!}{(l-m-1)!!}\,
  \frac{j_l(2\pi q)}{2\pi q} , \qquad l+m \text{ odd} .
  \label{eq:brief-odd}
\end{equation}
One spherical Bessel function, whatever the degree.

For $l+m$ even the factor is absent, and the quotient $P_{lm}(\sqrt{1-\rho^2})/\rho^m$ is a
polynomial in $\rho^2$ --- precisely because the $m$-th derivative is then an even function of
$x$.  That polynomial is expanded in Zernike radial polynomials of the same order,
\begin{equation}
  \begin{gathered}
    P_{lm}\bigl(\sqrt{1-\rho^2}\bigr) \;=\; \rho^m \sum_{k=0}^{(l-m)/2} c_k\,
    R_{m+2k}^{m}(\rho) , \\
    c_k \;=\; 2(m+2k+1)\int_0^1 \frac{P_{lm}\bigl(\sqrt{1-\rho^2}\bigr)}{\rho^m}\,
    R_{m+2k}^{m}(\rho)\,\rho\, d\rho .
  \end{gathered}
  \label{eq:brief-even-coefficients}
\end{equation}
The expansion is finite because both sides have the same degree and parity, and the Zernike
transform turns it term by term into ordinary Bessel functions,
\begin{equation}
  X_{lm}(q) \;=\; \sum_{k=0}^{(l-m)/2} c_k\,(-1)^k\,
  \frac{J_{m+2k+1}(2\pi q)}{2\pi q} , \qquad l+m \text{ even} .
  \label{eq:brief-even}
\end{equation}
So every mode of every degree is a finite Bessel sum, and a visibility at a single baseline
is a short sum over modes with no quadrature anywhere.

The expressions \eqref{eq:brief-odd} and \eqref{eq:brief-even}
were derived by \cite{dholakia2025} but the
expression \eqref{eq:brief-even-coefficients} giving the
coefficients $c_k$ as rational numbers is a refinement.


\begin{itemize}
  \item[$\checkmark$] \code{check_odd_closed_form} compares \eqref{eq:brief-odd} against the
    quadrature \eqref{eq:brief-radial-integral}, which is the definition.  This is the terminal
    check of the whole reduction, and it is the one to read.
    \\[1pt] \vtag{13}\quad \emph{Numerical to $10^{-9}$, $l \le 5$}
  \item[$\checkmark$] \code{check_odd_collapse_exact} obtains the same collapse coefficients a
    second way, symbolically.
    \\[1pt] \vtag{11}\quad \emph{Exact, $l \le 8$}
  \item[$\checkmark$] \code{check_k_identity} checks the $K$ identity the collapse rests on.
    \\[1pt] \vtag{18}\quad \emph{Numerical, $l \le 6$}
\end{itemize}

\begin{itemize}
  \item[$\checkmark$] \code{check_zernike_jacobi} relates the Zernike sum to the Jacobi form by
    expanding both as polynomials.
    \\[1pt] \vtag{8}\quad \emph{Exact, $n \le 8$}
  \item[$\checkmark$] \code{check_zernike_orthogonality} checks
    the Zernike orthogonality relation of section~\ref{sec:even}, in rational arithmetic for every pairing.
    \\[1pt] \vtag{9}\quad \emph{Exact, $n \le 8$}
  \item[$\checkmark$] \code{check_zernike_hankel} expands both sides of the Hankel identity in
    powers of $a$.
    \\[1pt] \vtag{10}\quad \emph{Exact, $n \le 6$, powers to $a^{18}$}
  \item[$\checkmark$] \code{check_zernike_projection} obtains the same coefficients by matching
    powers of $\rho$ downwards, a triangular system.
    \\[1pt] \vtag{12}\quad \emph{Exact rational, $l \le 6$}
  \item[$\checkmark$] \code{check_zernike_bessel_identity} checks the Zernike transform against
    the library's own $R_n^m$, by quadrature.
    \\[1pt] \vtag{22}\quad \emph{Numerical, $n \le 8$}
  \item[$\checkmark$] \code{check_even_closed_form} compares \eqref{eq:brief-even} against the
    same quadrature as V13.  This is the even family's terminal check: if it passes, the five
    lemmas above have been validated with it.
    \\[1pt] \vtag{14}\quad \emph{Numerical to $10^{-9}$, $l \le 6$}
\end{itemize}

\subsection{The deformed source}

A star deformed by its own rotation and by a companion has a level
surface given, to quadrupole order, by a potential of two amplitudes,
\begin{equation}
  \Phi(\mathbf{x}) \;=\; -\frac{1}{r}
  \;-\;\epsilon_\text{tide}\,r^2P_2(n\cdot\hat e_\text{tide})
  \;+\;\epsilon_\text{spin}\,r^2P_2(n\cdot\hat e_\text{spin})
  \;-\;\epsilon_\text{spin}\,r^2 ,
  \label{eq:brief-potential}
\end{equation}
whose surface is the level set $\Phi = -1$.  Writing the depth of the perturbation along a
direction as $E(n) = \epsilon_\text{tide}P_2(n\cdot\hat e_\text{tide})
- \epsilon_\text{spin}[P_2(n\cdot\hat e_\text{spin})-1]$, that surface is the solution of
$1/r + E(n)r^2 = 1$, which is exact and costs one root per sampled direction.

The brightness follows the gravity, $I \propto g^p$ with $g = |\nabla\Phi|$ on the surface and
$p = 1$ the bolometric von Zeipel law, multiplied by a linear limb factor
\begin{equation}
  1 - u_\text{limb}(1-\mu) , \qquad \mu = \hat n\cdot\hat e_\text{obs} ,
  \label{eq:brief-limb-law}
\end{equation}
which is linear in $\mu$ because $\mu$ is a single harmonic, and therefore costs one
coefficient at $m = 0$ rather than a product of expansions.

The shape and the brightness together make the map a function of position that no closed form
covers, so the scale of the calculation is set by how the surface is approximated.  Two
approximations are used.  The first replaces the sky map by its linear part: the shape is
captured by a single $2\times2$ positive-definite matrix, the stretch
\begin{equation}
  \Lambda \;=\; \sqrt{A A^{T}} , \qquad A = P\,(\text{frame}\cdot\text{radii}) ,
  \label{eq:brief-stretch}
\end{equation}
where $P$ projects onto $(\text{East},\text{North})$ and the radii are taken along the
principal directions.  Those directions are the eigenvectors of the quadratic form of
$E(n)$, and the radii along them are the exact level radii, so $\Lambda$ is computed from the
level surface rather than fitted.  The stretch acts on the visibility as a change of
frequency,
\begin{equation}
  V_\text{affine}(u,v) \;=\; V_\text{sph}\bigl(\Lambda^{T}(u,v)\bigr) ,
  \label{eq:brief-affine}
\end{equation}
with no flux factor: the determinant that stretching puts on the brightness cancels the one
it puts on the measure.  $\Lambda$ is best read as a preconditioner, since a different but
equally defensible stretch changes the coefficients of any series built on it without
changing the exact answer.

To first order in both amplitudes the pattern is a constant plus two quadrupoles, and with
the limb factor carried through it becomes
\begin{equation}
  I(\hat n) \;=\; \bigl(1 - u_\text{limb}(1-\mu)\bigr)\Bigl[\,1 - 4\epsilon_\text{spin}
  - 4\epsilon_\text{tide}P_2(\hat n\cdot\hat t)
  + 4\epsilon_\text{spin}P_2(\hat n\cdot\hat s)\Bigr] ,
  \label{eq:brief-first-order}
\end{equation}
whose coefficients are, to leading order in $u_\text{limb}\epsilon$,
\begin{equation}
  \begin{aligned}
    a_{00} &= (1-u_\text{limb})\,(1-4\epsilon_\text{spin}) , \\
    a_{10} &= \frac{u_\text{limb}}{\sqrt3} , \qquad
    a_{2m} = -\frac{4}{5}\Bigl[\epsilon_\text{tide}\,\overline{Y_{2m}(\hat t)}
    - \epsilon_\text{spin}\,\overline{Y_{2m}(\hat s)}\Bigr] .
  \end{aligned}
  \label{eq:brief-first-order-coefficients}
\end{equation}
A constant, a dipole from the limb, and two quadrupoles from the deformation, with no Gaunt
integral anywhere: the limb is a dipole because $\mu$ is a single harmonic, and the
quadrupoles come from \eqref{eq:brief-zonal} and \eqref{eq:brief-y2m}.

The second approximation does not assume the coefficients at all.  The surface is rendered,
sampled on the level surface with the limb factor applied, and coefficients fitted by least
squares in the image plane, which is the norm the data live in.  The fit solves a normal
equation whose matrix is a Gram matrix of the basis alone, and whose right-hand side is a sum
over the samples,
\begin{equation}
  b_{lm} \;=\; \sum_s w_s\, I_s\, \overline{\Phi_{lm}\bigl(\Lambda^{-1}x_s\bigr)} ,
  \label{eq:brief-fitted-rhs}
\end{equation}
with $\Phi_{lm}$ the projected harmonics and $w_s$ the projected area each sample carries.
This row has no assumption beyond the stretch, so its error is a truncation at the chosen
degree, which the user controls; the first-order row's error is instead the linearisation of
the brightness.  Comparing the two is what the interactive models are for.




\begin{itemize}
  \item[$\checkmark$] \code{check_deformation_first_order} checks the first-order relations
    between $\delta r$, $\delta g$ and the potential.
    \\[1pt] \vtag{19}\quad \emph{Algebraic}
  \item[$\checkmark$] \code{check_level_surface} evaluates $\Phi$ on the returned surface and
    checks $\nabla\Phi$ against a central difference.  A sign error in the potential fails it
    immediately.
    \\[1pt] \vtag{20}\quad \emph{Numerical}
\end{itemize}

\begin{itemize}
  \item[$\checkmark$] \code{check_first_order_coefficients} compares the closed-form coefficients
    \eqref{eq:brief-first-order-coefficients} against the quadrature that projects the same
    pattern.  It covers the limb factor as well, which is why both displays share it.
    \\[1pt] \vtag{21}\quad \emph{Numerical}
  \item[$\checkmark$] \code{check_first_order_frame} compares the coefficients computed with
    body-frame axes against those rotated into the observer frame.  Both routes share the
    $l = 2$ corollary, so a wrong constant would pass; a wrong frame rule would not.
    \\[1pt] \vtag{24}\quad \emph{Algebraic}
\end{itemize}

\begin{itemize}
  \item[$\checkmark$] \code{check_fitted_round_trip} checks that the fit returns a synthetic
    image's own coefficients, on samples that need no quadrature.
    \\[1pt] \vtag{26}\quad \emph{Algebraic}
  \item[$\checkmark$] \code{check_disk_gram} checks the Gram matrix's analytic entry, its block
    structure and its midpoint rule.
    \\[1pt] \vtag{27}\quad \emph{Numerical}
\end{itemize}

\subsection{What is new here}

Sections 1--3 are standard spherical-harmonic material, and the reduction of
\eqref{eq:brief-mode}--\eqref{eq:brief-radial-integral} follows published work.  The
closed forms of \eqref{eq:brief-odd} and \eqref{eq:brief-even} use the same parity split
as Dholakia \& Pope (2025), but give each mode as an explicit Bessel sum rather than as a
basis-change matrix.  The deformation of
\eqref{eq:brief-potential}--\eqref{eq:brief-first-order-coefficients}, with the limb
carried in the coefficients, and the fitted row of \eqref{eq:brief-fitted-rhs}, are the new
material, and the two interactive models that compare them are the second.

